% Cross-architecture comparison tables (Ada sm_89 vs A100 sm_80).
% AUTO-GENERATED from experiments/results/<gpu>/*.csv and ViperBench/results/profile.*.csv
% by experiments/results/cross_arch/ generator -- numbers are NOT hand-typed.
% Answers R1-Q5 / R3-Q2 (do conclusions generalize to a data-center GPU?).

\begin{table}[t]
  \centering
  \caption{\textbf{Cross-architecture generalization.} Median library efficiency $E_\text{lib}$ (\%) per
    kernel category for each DSL on the workstation RTX 4000 Ada (sm\_89) and the data-center
    A100-PCIE (sm\_80), untuned defaults. The qualitative profile is preserved across architectures:
    GEMM and element-wise competitive, convolution and TileLang normalization severely behind.}
  \label{tab:xarch:summary}
  \begin{tabular}{lcccc}
    \toprule
    & \multicolumn{2}{c}{Triton} & \multicolumn{2}{c}{TileLang} \\
    \cmidrule(lr){2-3}\cmidrule(lr){4-5}
    Category & Ada & A100 & Ada & A100 \\
    \midrule
    GEMM & 55.1\% & 75.5\% & 55.2\% & 60.3\% \\
    Convolution & 42.8\% & 38.0\% & 8.8\% & 11.1\% \\
    Normalization & 123.5\% & 137.3\% & 1.5\% & 0.8\% \\
    Element-wise/Reduction & 97.3\% & 73.0\% & 67.0\% & 55.3\% \\
    \bottomrule
  \end{tabular}
\end{table}

\begin{table}[t]
  \centering
  \caption{\textbf{Root-cause taxonomy reproduces across architectures.} Each corrected root cause
    measured on both GPUs with the identical portable harness. All four reproduce, confirming they
    are properties of the DSL/compiler rather than of the Ada hardware (cf.\ \cref{tab:xarch:summary}).}
  \label{tab:xarch:rootcause}
  \resizebox{\columnwidth}{!}{
  \begin{tabular}{llll}
    \toprule
    Root cause & Ada (sm\_89) & A100 (sm\_80) & Reproduces \\
    \midrule
    RC0 TileLang LayerNorm anomaly ($8192^2$) & $E_\text{lib}=0.32\%$ & $E_\text{lib}=0.09\%$ & yes (memory-latency bound) \\
    RC3 Triton conv register spill (1$\times$1--7$\times$7) & $n_\text{spills}=0$ & $n_\text{spills}=0$ & yes (no spill; occupancy) \\
    RC4 Winograd upper-bound contribution & $2.3\%$ & $2.0\%$ & yes ($\sim$2--3\%, not primary) \\
    FP32 \texttt{T.gemm} silent TF32 lowering & 3013$\times$ rel.\ err & 633$\times$ rel.\ err & yes (matches TF32 class) \\
    \bottomrule
  \end{tabular}}
\end{table}

\begin{table}[t]
  \centering
  \caption{\textbf{GEMM autotuning recovery generalizes (RC2).} Library efficiency of the Triton
    matmul before (plain, heuristic) and after (expanded autotune search) on both GPUs. The
    \S5-vs-\S7.3 gap is a search-space artifact on both architectures, not shape- or hardware-specific.}
  \label{tab:xarch:autotune}
  \begin{tabular}{lcccc}
    \toprule
    & \multicolumn{2}{c}{Triton plain} & \multicolumn{2}{c}{Triton autotuned} \\
    \cmidrule(lr){2-3}\cmidrule(lr){4-5}
    Shape & Ada & A100 & Ada & A100 \\
    \midrule
    $4096^2$ & 83.0\% & 57.7\% & 102.2\% & 110.1\% \\
    $16384^2$ & 32.2\% & 38.3\% & 98.5\% & 87.5\% \\
    \bottomrule
  \end{tabular}
\end{table}

\begin{table}[t]
  \centering
  \caption{\textbf{Convolution filter sweep across architectures} (large shape $32{\times}256{\times}128^2$,
    groups$=$1, stride$=$1). $E_\text{lib}$ vs cuDNN for each DSL, and the measured Triton register-spill
    count. The gap and the absence of register spilling both hold across architectures.}
  \label{tab:xarch:conv}
  \begin{tabular}{lccccc}
    \toprule
    & \multicolumn{2}{c}{Triton $E_\text{lib}$} & \multicolumn{2}{c}{TileLang $E_\text{lib}$} & Triton \\
    \cmidrule(lr){2-3}\cmidrule(lr){4-5}
    Filter & Ada & A100 & Ada & A100 & $n_\text{spills}$ \\
    \midrule
    $1\times1$ & 23.3\% & 37.0\% & 11.5\% & 10.3\% & 0 \\
    $3\times3$ & 36.5\% & 24.4\% & 9.4\% & 15.3\% & 0 \\
    $5\times5$ & 15.8\% & 16.9\% & --- & 12.9\% & 0 \\
    $7\times7$ & 12.9\% & 12.5\% & --- & --- & 0 \\
    \bottomrule
  \end{tabular}
\end{table}
